\chapter{一般群的 Langlands 纲领 I}

% ===== PDF p.088 =====
\noindent\textbf{接下来的目标}
\begin{itemize}
\item 对 $\GL_n$、乃至 $\Q$ 上一般的约化群 $G$，给出（部分）Langlands 纲领的或多或少精确的陈述。
\end{itemize}

\medskip
\noindent\textbf{$\GL_{n,F}$ 的 Langlands（–Clozel–Fontaine–Mazur）互反}
\begin{itemize}
\item 设 $F$ 为任意数域。
\item 固定某个素数 $\ell$ 与同构 $\C\cong\Ql$。
\item Langlands（–Clozel–Fontaine–Mazur）互反猜想断言：对任意 $n\ge1$，存在一个（唯一的）双射，介于
\begin{itemize}
\item[(i)] $\GL_n(\A_F)$ 的 $L$-代数尖点自守表示之集，
\item[(ii)] 不可约连续表示 $\Gal(\overline F/F)\to\GL_n(\Ql)$ 的同构类之集，这些表示几乎处处非分歧，且在整除 $\ell$ 的位处为 de Rham 表示，
\end{itemize}
之间，使得该双射把 Satake 参数与 Frobenius 元素的特征值相匹配。\znote{原文此处作 “Langland(-Clozel-Fontaine-Mazur)”，应为 “Langlands”，漏了一个 s；译文按 Langlands 处理。}
\end{itemize}

\medskip
\noindent\textbf{评论}
\begin{itemize}
\item 这只是 $G=\Res_{F/\Q}\GL_{n,F}$ 的 Langlands 纲领的一部分。
\item 它既没有描述 $\GL_{n,F}$ 的全部自守表示，也没有描述全部的 Galois 表示。
\item 对 $F=\Q$、$n\ge2$ 未知。例如：$\GL_2(\A)$ 存在 $L$-代数尖点自守表示（与 Maass 形式相伴），它们应当对应于 $\Gal(\overline\Q/\Q)$ 的偶的、不可约的二维 $\ell$-进表示（像有限）。
\item Langlands–Tunnell：若像可解，则该 Galois 表示的自守性已知，参见 \texttt{[GH19, Theorem 13.4.5.]}。这并不覆盖像为 $\SL_2(\F_5)$ 的、几乎处处非分歧的二维不可约 Galois 表示。
\item 唯一性可由双方的（强）重数一定理得到。
\item 今日：解释「$L$-代数」是什么意思。
\end{itemize}

\medskip
\noindent\textbf{一般群的自守形式}
\begin{itemize}
\item $G$ 为 $\Q$ 上任意的约化群。
\item 我们将为 $L^2([G])$ 引入一个方便的（= 更代数的）替代物。
\item 回忆第 3 节中，对 $G=\GL_2=\GL_{2,\Q}$、$k\in\Z$ 的嵌入
\[
\Phi\colon M_k\longrightarrow C^\infty(\GL_2(\A)),\qquad f\longmapsto\phi_f.
\]
\item 回忆 $\Phi(M_k)$ 与 $\Phi(S_k)$ 的像的描述，即：
\begin{itemize}
\item 对 $k\in\Z$，元素 $\phi\in C^\infty(\GL_2(\A))$ 落在 $\Phi\colon M_k\to C^\infty(\GL_2(\A))$ 的像中，当且仅当
\begin{itemize}
\item 对 $(g,g_\infty)\in\GL_2(\A)$、$z\in\C^\times\subseteq\GL_2(\R)$，有 $\phi(g,g_\infty z)=z^{-k}\phi(g,g_\infty)$；
\item $\phi$ 是适度增长的（这蕴涵负 Fourier 系数为零）；
\item 对每个 $g\in\GL_2(\Af)$ 有 $Y*\phi(g,-)=0$，
% ===== PDF p.089 =====
其中 $Y\in\g_\C=\mathfrak{gl}_2(\R)_\C$ 是适当构造的自然元素（这蕴涵全纯性）；
\item 对一切 $\gamma\in\GL_2(\Q)$ 有 $\phi(\gamma g)=\phi(g)$。
\end{itemize}
\item $f$ 还是尖点的\idx{尖点的}，当且仅当对一切含于真抛物子群 $P\subseteq G$（定义在 $\Q$ 上）的幺幂根基 $N$，有
\[
\int_{N(\Q)\backslash N(\A)}\phi_f(ng)\,\mathrm{d}n=0.
\]
\end{itemize}
\item 固定 $G(\R)$ 的一个极大紧子群 $K_\infty$（通常不连通），例如 $\mathrm{O}_2(\R)=\SO_2(\R)\cup s\SO_2(\R)\subseteq\GL_2(\R)$，其中
\[
s=\begin{pmatrix}-1&0\\0&1\end{pmatrix}.
\]
\end{itemize}

\medskip
\noindent\textbf{定义 11.1}（参见 \texttt{[GH19, Definition 6.5]}）\phantomsection\label{def:11.1}\quad 设 $G$ 为 $\Q$ 上的约化群。$G$ 的一个（阿代尔）\textbf{自守形式}是光滑函数 $\phi\colon G(\A)\to\C$，使得
\begin{itemize}
\item[1)] $\phi$ 是适度增长的（第 4 节已定义；本讲不需要）；
\item[2)] $\phi$ 是右 $K_\infty$-有限的，即函数 $g\mapsto\phi(gk)$（$k\in K_\infty\subseteq G(\A)$）张成有限维向量空间；
\item[3)] $\phi$ 被 $\cZ(\g_\C)$ 中一个有限余维的理想杀死，其中 $\g_\C$ 记 $G(\R)$ 的（复化）Lie 代数，$\cZ(\g_\C)$ 记 $\g_\C$ 的泛包络代数 $U(\g_\C)$ 的中心；
\item[4)] $\phi$ 是左 $G(\Q)$-不变的，即对一切 $\gamma\in G(\Q)$ 有 $\phi(\gamma g)=\phi(g)$。
\end{itemize}

\medskip
\noindent\textbf{$\GL_2$ 的 Casimir 元}
\begin{itemize}
\item 我们来验证：对 $f\in M_k$，函数 $\phi_f\in C^\infty(\GL_2(\A))$ 确实满足 3)。
\item 也就是说，元素 $Y$ 并不在 $\cZ(\mathfrak{gl}_{2,\C})$ 中。
\item 置
\[
H:=\begin{pmatrix}0&-i\\ i&0\end{pmatrix},\qquad
X:=\frac12\begin{pmatrix}-1&-i\\ -i&1\end{pmatrix},\qquad
Y:=\frac12\begin{pmatrix}-1&i\\ i&1\end{pmatrix},\qquad
Z:=\begin{pmatrix}1&0\\ 0&1\end{pmatrix}.
\]
\item 于是
\begin{itemize}
\item[1)] $H*\phi_f=-k\phi_f$，因为对 $z\in\C^\times$ 有 $z\cdot\phi_f=z^{-k}\phi_f$；
\item[2)] $Y*\phi_f=0$。
\end{itemize}
\item $[X,Y]=XY-YX=H$（因为 $H,X,Y$ 构成 $\mathfrak{sl}_2$-三元组）。
\item 定义 Casimir 算子
\begin{equation}\tag{1}\label{eq:11.1}
\Delta:=\tfrac14\bigl(H^2+2XY+2YX\bigr).
\end{equation}
则 $\Delta\in\cZ(\mathfrak{gl}_{2,\C})$。
\item 事实上，$\cZ(\mathfrak{gl}_{2,\C})$ 作为 $\C$-代数由 $\Delta$ 与 $Z\in\mathfrak{gl}_{2,\C}$ 生成（这可由 Harish-Chandra 同构推出，参见 \texttt{[GH19, Theorem 4.6.1]}）。
\item 计算（利用 $Y*\phi_f=0$）得
\begin{equation}\tag{2}\label{eq:11.2}
\Delta*\phi_f=\tfrac14\bigl(H^2+4YX+2H\bigr)*\phi_f=\tfrac14\bigl(H^2+2H\bigr)*\phi_f=\tfrac14(k^2-2k)\phi_f,
\end{equation}
于是 $\phi_f$ 确实被 $\cZ(\mathfrak{gl}_{2,\C})$ 中有限余维的理想 $\bigl\langle\Delta-\tfrac14(k^2-2k),\,Z+k\bigr\rangle$ 杀死。\znote{原文此处作 $\Delta*\phi_f=\frac14(H^2-2XY+2YX)*\phi_f=\frac14(H^2-2H)*\phi_f=\frac14(k^2-2k)\phi_f$：中间式与 $\Delta$ 的定义不符，且由 $H*\phi_f=-k\phi_f$，该式会给出 $\frac14(k^2+2k)\phi_f$，与末式矛盾。按 $[X,Y]=H$（即 $XY=YX+H$）有 $2XY+2YX=4YX+2H$；再由 $Y*\phi_f=0$ 得 $YX*\phi_f=0$，故 $\Delta*\phi_f=\frac14(H^2+2H)*\phi_f=\frac14(k^2-2k)\phi_f$。译文按更正后的计算写出。[来源：底本 p.~89；见疑点清单 p.89 条]}
\end{itemize}

% ===== PDF p.090 =====
\medskip
\noindent\textbf{关于自守形式的补充}
\begin{itemize}
\item $G$ 为 $\Q$ 上任意的约化群。
\item 记 $\cA(G)$ 为自守形式构成的空间，$\cA([G])$ 为其中在 $A_G$（经 $G(\A)$ 上的左平移或右平移作用）下不变的自守形式构成的子空间。\znoteplain{记号 $[G]$ 沿用第 2 讲：$[G]=G(\Q)A_G\backslash G(\A)$。}
\item 遗憾的是，空间 $\cA([G])$ 与 $\cA(G)$ 在 $G(\R)$ 下并不稳定，因为 $K_\infty$-有限性未必被保持。
\end{itemize}

\medskip
\noindent\textbf{$K_\infty$-有限性不被保持的例子（据 Andreas Mihatsch）}
\begin{itemize}
\item 取 $G=\GL_2(\R)$、$K:=K_\infty^\circ=\SO_2(\R)\cong S^1$，并取 $V$ 为 $G$ 上函数空间中的任一表示，只要它包含 $C_c^\infty(G)$。我们给出一个例子，说明 $K$-有限性在右平移作用下不被保持。
\item 更精确地说：存在 $\phi\in C_c^\infty(G)$ 与 $h\in G$，使得 $\phi$ 是 $K$-有限的，但不是 $K_h:=h^{-1}Kh$-有限的。
\item 为此：取 $h\in G$ 使 $K_h\cap K=\{\pm1\}$。
\item 再取非零的 $\phi\in C_c^\infty(G)$，它在 $K$ 下不变，且在 $K_h$ 的某个开子集上 $\phi\equiv0$（例如：把 $G/K$ 上的函数拉回以构造 $\phi$；注意 $K_h$ 在 $G/K$ 中的像维数 $>0$）。
\item $K_h$ 上任何非零的 $K_h$-有限函数都不能在开子集上为零。事实上，$K_h\cong S^1$ 上的 $K_h$-有限函数恰为 Laurent 多项式的限制。
\end{itemize}

\medskip
\noindent\textbf{关于自守形式的补充（续）}
\begin{itemize}
\item 显然，$G(\Af)$ 经 $\phi\mapsto\bigl(g\mapsto\phi(gh)\bigr)$（$h\in G(\Af)$）作用在 $\cA(G)$ 上。
\item 类似地，$K_\infty$、乃至 $A_GK_\infty$，作用在 $\cA(G)$ 上。
\item 不那么显然、但为真的是（参见 \texttt{[GH19, Proposition 4.4.2]}）：$\g_\C$ 作用在 $\cA(G)$ 上。$\Rightarrow$ $U(\g_\C)$、特别是 $\cZ(\g_\C)\subseteq U(\g_\C)$，经微分算子作用在 $\cA(G)$ 上。
\item $\g_\C$-作用与 $K_\infty$-作用使 $\cA(G)$ 成为下述意义下的 $(\g_\C,K_\infty)$-模。
\end{itemize}

\medskip
\noindent\textbf{定义 11.2}（参见 \texttt{[GH19, Definition 4.5]}）\phantomsection\label{def:11.2}\quad 一个 $(\g_\C,K_\infty)$-模是一个 $\C$-向量空间 $V$，它同时是 $\g_\C$-模与 $K_\infty$-模，且满足
\begin{itemize}
\item[1)] $V$ 是有限维 $K_\infty$-稳定子空间的并；
\item[2)] 对 $X\in\mathfrak k:=\Lie(K_\infty)_\C$ 与 $\phi\in V$，有
\[
X*\phi=\frac{\partial}{\partial t}\bigl(\exp(tX)\phi\bigr)\Big|_{t=0};
\]
\item[3)] 对 $h\in K_\infty$、$X\in\g_\C$ 与 $\phi\in V$，有
\[
h\bigl(X*(h^{-1}\phi)\bigr)=\bigl(\mathrm{Ad}(h)(X)\bigr)*\phi.
\]
\end{itemize}

\medskip
\noindent\textbf{评论}
\begin{itemize}
\item 我们对 $V$ 不要求任何拓扑，因此 $(\g_\C,K_\infty)$-模比 $G(\R)$-表示代数得多，也就更容易处理。
\item 2) 中的极限是对 $V$ 中某个包含 $\phi$ 的有限维 $K_\infty$-稳定子空间的自然拓扑取的。
\item 对 $G(\R)$ 在某个 Hilbert 空间（视为 $\C$-向量空间）上的每个表示，光滑且 $K_\infty$-有限的向量构成的空间自然地是 $(\g_\C,K_\infty)$-模，参见 \texttt{[GH19, Proposition 4.4.2]}；而对酉表示不会丢失任何信息，参见 \texttt{[GH19, Theorem 4.4.4]}，即：$G(\R)$ 的不可约酉表示的等类单射到不可约 $(\g_\C,K_\infty)$-模的等类中。
\item 此外，如此得到的 $(\g_\C,K_\infty)$-模是容许的，即：对 $K_\infty$ 的每个（有限维）不可约表示 $\sigma$，其 $\sigma$-等型分量是有限维的，参见 \texttt{[GH19, Theorem 4.4.1]}。
\item Schur 引理的一个变体：在每个不可约容许 $(\g_\C,K_\infty)$-模 $\pi$ 上，$U(\g_\C)$ 的中心 $\cZ(\g_\C)$ 经某个 $\C$-代数同态 $\omega_\pi\colon\cZ(\g_\C)\to\C$ 作用；$\omega_\pi$ 称为 $\pi$ 的\textbf{无穷小特征}（infinitesimal character）。
\item Harish-Chandra 同构（参见 \texttt{[GH19, Theorem 4.6.1]}）给出：对每个 Cartan 子代数 $\mathfrak t\subseteq\g$ 与相应的 Weyl 群 $W$，有 $\cZ(\g_\C)\cong U(\mathfrak t_\C)^W$。注意 $U(\mathfrak t_\C)\cong\Sym^\bullet\mathfrak t_\C$，因为 $\mathfrak t_\C$ 是 abelian Lie 代数。
\item 由于 $W$ 有限，$\C$-代数同态 $\cZ(\g_\C)\cong U(\mathfrak t_\C)^W\to\C$ 与 $\C$-代数同态 $U(\mathfrak t_\C)\cong\Sym^\bullet\mathfrak t_\C\to\C$ 的 $W$-轨道一一对应。此外
\[
\Hom_{\C\text{-}\mathrm{alg}}\bigl(U(\mathfrak t_\C),\C\bigr)\cong\Hom_{\C\text{-}\mathrm{lin}}\bigl(\mathfrak t_\C,\C\bigr)=:\mathfrak t_\C^\vee,
\]
且 $U(\mathfrak t_\C)^W\hookrightarrow U(\mathfrak t_\C)$ 是有限扩张。
\item 于是无穷小特征定义了「权空间」$\mathfrak t_\C^\vee$ 中一个典范的 $W$-轨道。
\item 现在设对某个极大环面 $T\subseteq G_\R$ 有 $\mathfrak t=\Lie(T)$。则 $W$-等变地有 $X^*(T_\C)\subseteq\mathfrak t_\C^\vee$。
% ===== PDF p.091 =====
\item 不可约容许 $(\g_\C,K_\infty)$-模 $\pi$ 称为\textbf{$L$-代数的}，若与无穷小特征 $\omega_\pi$ 相应的 $\mathfrak t_\C^\vee$ 中的 $W$-轨道落在有限自由 $\Z$-模 $X^*(T_\C)\subseteq\mathfrak t_\C^\vee$ 中，参见 \texttt{[BG10, Definition 2.3.1]}、\texttt{[Tho, Definition 93]}。
\end{itemize}

\medskip
\noindent\textbf{由模形式生成的 $(\mathfrak{gl}_{2,\C},\mathrm{O}_2(\R))$-模}
\begin{itemize}
\item 设 $f\in M_k$。我们来分析由 $\phi:=\phi_f$ 生成的 $(\mathfrak{gl}_{2,\C},\mathrm{O}_2(\R))$-模 $V\subseteq\cA(\GL_2)$。
\item 回忆 $Y*\phi=0$，且对 $z\in\C^\times$ 有 $z\cdot\phi=z^{-k}\phi$。
\item 不难看出：存在唯一的（在同构意义下）由这样的元素 $\phi$ 生成的 $(\mathfrak{gl}_{2,\C},\mathrm{O}_2(\R))$-模 $D'_{k-1}$（即 \texttt{[Del73, Section 2.1]} 中 $D_{k-1}$ 的对偶），并且 $D'_{k-1}\cong V$ 不可约。
\item 特别地，它只依赖于 $k$，不依赖于 $f$。
\item 正如之前某处提到的，我们得到
\[
M_k\cong\Hom_{(\mathfrak{gl}_{2,\C},\mathrm{O}_2(\R))}\bigl(D'_{k-1},\cA(\GL_2)\bigr)
\]
（注：\texttt{[Del73, Scholie 2.1.3]} 中的空间 $H(G_\A)$ 与我们的 $\cA(\GL_2)$ 只在 $\GL_2(\A)$ 上取逆的意义下一致）。
% ===== PDF p.092 =====
\item 我们算得 $\Delta*\phi=\tfrac14(k^2-2k)\phi$、$Z*\phi=-k\phi$，这确定了无穷小特征 $\omega_{D'_{k-1}}$。
\item 在 Harish-Chandra 同构下，与 $\omega_{D'_{k-1}}$ 相应的 $\C$-代数同态 $U(\mathfrak t_\C)\cong\C[Z,H]\to\C$ 的 $W$-轨道由同态
\[
Z\longmapsto-k,\qquad H\longmapsto\pm(k-1)
\]
给出。
\item 格 $X^*(T)\subseteq\langle Z,H\rangle^\vee_\C$ 由满足 $a,b\in\Z$ 且 $a\equiv b\pmod2$ 的 $\C$-线性映射 $Z\mapsto a$、$H\mapsto b$ 给出。
\item 因此 $D'_{k-1}$ 不是 $L$-代数的（它在 \texttt{[BG10]} 的意义下是 $\C$-代数的）。\znoteplain{这里 $L$-代数性要求该 $W$-轨道落在 $X^*(T)$ 中，即要求 $Z\mapsto-k$ 与 $H\mapsto\pm(k-1)$ 这两个坐标同奇偶；但 $-k$ 与 $k-1$ 的奇偶性总是相反，故 $D'_{k-1}$ 从不 $L$-代数。扭上 $|\det|^{1/2+\alpha}$ 之后，$Z$ 的坐标变为 $1-k+2\alpha$，它与 $k-1$ 同奇偶，故扭后确是 $L$-代数的。另外注意这里用的是特征格 $X^*(T)$（$W$-等变地含于权空间 $\mathfrak t^\vee_\C$），而非余特征格 $X_*(T)$。}
\item 然而，对任意 $a\in\Z$，$D'_{k-1}\otimes_\C|\det|^{1/2+a}$ 是 $L$-代数的，因为（与 $|\det|^{1/2}$ 相伴的 $(\mathfrak{gl}_{2,\C},\mathrm{O}_2(\R))$-模）的无穷小特征对应于 $Z\mapsto1$、$H\mapsto0$。
\end{itemize}

\medskip
\noindent\textbf{要点}
\begin{itemize}
\item 对任意 $G$ 都存在（阿代尔）自守形式的空间 $\cA(G)$，当 $G=\GL_2$ 时它自然地包含诸 $\phi_f$（$f\in M_k$）。
\item 空间 $\cA(G)$ 带有 $(\g_\C,K_\infty)\times G(\Af)$-作用，是带 $G(\A)$-作用的 $L^2([G])$ 的一个更代数的替代物。
\end{itemize}

\medskip
\noindent\textbf{自守表示的一个新概念}
\begin{itemize}
\item 我们用 $\cA(G)$ 代替 $L^2([G])$ 来重新定义自守表示。
\item 即：$G$ 的自守表示是 $\cA(G)$ 的不可约 $G(\Af)\times(\g_\C,K_\infty)$-子商，参见 \texttt{[GH19, Definition 6.8]}。
\item 此前的自守表示今后称为「$L^2$-自守表示」。
\item 两种概念之间的关系参见 \texttt{[GH19, Section 6.5]}。
\end{itemize}

\medskip
\noindent\textbf{Flath 定理（参见 \texttt{[GH19, Theorem 5.7.1]}）}
\begin{itemize}
\item 对几乎所有素数 $p$，约化群 $G_{\Q_p}:=G\times_{\Spec(\Q)}\Spec(\Q_p)$ 是「非分歧的」，即它可以延拓为约化群概形 $G_p\to\Spec(\Z_p)$。
\item 等价地，$G_{\Q_p}$ 在一个非分歧扩张上是准分裂的（即含一个定义在 $\Q_p$ 上的 Borel 子群）并且是分裂的（即含一个极大分裂环面）。
\item 于是
\begin{equation}\tag{3}\label{eq:11.3}
G(\A)=\sideset{}{'}\prod_{p}\bigl(G(\Q_p),G_p(\Z_p)\bigr)\times G(\R),
\end{equation}
其中 $G_p(\Z_p)\subseteq G(\Q_p)$ 是紧开拓子群（参见 \texttt{[GH19, Proposition 2.3.1]}），称为\textbf{超特殊的}。
% ===== PDF p.093 =====
\item Flath 定理（参见 \texttt{[GH19, Section 5.7]} 与 \texttt{[Fla79]}）断言：不可约容许 $G(\Af)\times(\g_\C,K_\infty)$-模 $\pi$（唯一地）分解为「限制张量积」
\begin{equation}\tag{4}\label{eq:11.4}
\pi\cong\sideset{}{'}\bigotimes_{p\ \text{素数}}\pi_p\otimes\pi_\infty,
\end{equation}
其中 $\pi_p$ 是不可约光滑（甚至是容许的）$G_{\Q_p}$-表示，$\pi_\infty$ 是不可约容许 $(\g_\C,K_\infty)$-模。
\item 此外，对几乎所有素数 $p$，表示 $\pi_p$ 非分歧，即群 $G_{\Q_p}$ 非分歧且 $\pi_p^{G_p(\Z_p)}\ne0$，其中 $G_p$ 是 $G_{\Q_p}$ 的约化模型。
\item 重要（参见 \texttt{[GH19, Theorem 7.5.1]}）：$\cH(G(\Q_p),G_p(\Z_p))$ 仍然是交换的！（$\Rightarrow$ 若 $\pi_p^{G_p(\Z_p)}\ne0$，则其维数为 $1$。）
\end{itemize}

\medskip
\noindent\textbf{要点}
\begin{itemize}
\item 于是不可约容许 $G(\Af)\times(\g_\C,K_\infty)$-模是一族「局部数据」。
\item 而「是自守的」在这些局部数据之间加上了很强的关系。
\item 对几乎所有素数 $p$ 得到同态
\[
\chi_p\colon\cH(G(\Q_p),G_p(\Z_p))\longrightarrow\C,
\]
与 $\GL_2$ 的情形类似；$\Rightarrow$ 这将给出「Hecke 特征系统」的类比。
\item 对 $\pi$ 置 $\pi_f:=\sideset{}{'}\bigotimes_p\pi_p$，于是 $\pi\cong\pi_f\otimes\pi_\infty$。
\end{itemize}

\medskip
\noindent\textbf{定义 11.3}（$L$-代数自守表示）\phantomsection\label{def:11.3}\quad 自守表示 $\pi\cong\pi_f\otimes\pi_\infty$ 称为\textbf{$L$-代数的}，若不可约容许 $(\g_\C,K_\infty)$-模 $\pi_\infty$ 是 $L$-代数的，参见 \texttt{[BG10, Definition 3.1.1]}。

\medskip
\noindent\textbf{例}
\begin{itemize}
\item 由模形式生成的自守表示 $\pi\subseteq\cA(\GL_2)$ 不是 $L$-代数的，但对任意 $a\in\Z$，其扭 $\pi\otimes_\C|\det|^{1/2+a}_{\text{adèlic}}$ 是 $L$-代数的。
\item 连续特征标 $\psi\colon\Q^\times\backslash\A^\times\to\C$ 是 $L$-代数的，当且仅当对某个 $k\in\Z$ 有 $\psi=\chi|\cdot|^k_{\text{adèlic}}$，其中 $|\cdot|_{\text{adèlic}}$ 是阿代尔范数，$\chi\colon\Q^\times\backslash\A^\times\to\C$ 是有限阶特征标。
\end{itemize}
